%%%%% Section 6. Demonstration with real data %%%%%
\section{Demonstration}\label{sec:demonstration}

Sections~\ref{sec:framework}--\ref{sec:computation} built the library
in concept, in mathematics, and in computation. What remains is to
build one. This section constructs the OT-FB from two actual
collections: a music collection, whose content data come from sound
alone, and a collection of children's books, which carries two kinds
of content data and so receives two bundles of rabbit holes at once.
We first state what the construction asks of the content data, then
build the two libraries in turn.


%%% Subsection 6.1 %%%
\subsection{What the Construction Asks of the Content
    Data}\label{subsec:content_data}

The OT-FB inherits its geometry from the content space. Everything the
construction of Sections~\ref{sec:mathematics}
and~\ref{sec:computation} rests on is measured there: the proximity
assumption (A2) and the transport cost are both stated in the
Euclidean distance of $\mathcal{X}$---the general-purpose inner
product that Section~\ref{sec:framework} asked the content data to
carry. Which works a rabbit hole joins is therefore decided by what
counts as close in $\mathcal{X}$.

The content space must accordingly be \emph{unfolded}: wherever works
are kindred, Euclidean closeness must report it, and wherever works
lie close, kinship must be what brought them there. The demand falls
on neighborhoods, not on the space at large; the space may curve as
it pleases, so long as nearness means kinship. Raw signals do not
satisfy this. Two kindred performances may lie far apart as
waveforms, and two kindred books as word counts, while accidents of
encoding sit side by side; a bundle built on such a space would dig
its holes along those accidents.

Deep embeddings supply spaces of the required kind. Foundation
encoders are trained on large corpora precisely so that semantic
closeness becomes geometric closeness, and their embedding spaces are
unfolded in the sense we need. For sound, we use CLAP~\cite{Wu2023},
a joint audio--text embedding: the music library of
Section~\ref{subsec:music_library} is built from sound alone, with no
words about the works entering the content data.

When the content data are words about the works, an LLM text
embedding plays the same role. The words need not come from a
controlled vocabulary: a ready-made description---an abstract, a
synopsis---serves as it stands, and where none exists, one is
composed from the tags users have freely given---``\textit{this work
carries the following tags: $t_1, t_2, \ldots$}''---and embedded
whole. Nor is a word any longer a label:
\textit{chill} and \textit{relaxed} lie close together in the
embedding, as they do in use. This is, in the end, the visitor's own
method---``\textit{quiet, but with feeling in it}'' was already a
description, offered in whatever words came; an LLM embedding
receives such words and returns a place.

Here one further choice opens: which words are embedded decides what
the bundle carries. Embed the adjectives given to a work, and its
atmosphere enters the content space; embed the nouns, and its themes
do. Section~\ref{subsec:book_library} uses exactly this freedom to
give one collection two kinds of rabbit holes.


%%% Subsection 6.2 %%%
\subsection{A Music Library Built from Sound
    Alone}\label{subsec:music_library}

% MTG-Jamendo (5 genres, exclusive tagging) + LAION-CLAP.
% Content data: audio only; norm normalization + PCA -> X.
% Genre only is given for training; mood/theme tags are held out
% and meet V afterward.
% TODO: 実験結果待ち。パラグラフ設計は結果より先に行う。


%%% Subsection 6.3 %%%
\subsection{A Book Library with Two Kinds of Rabbit
    Holes}\label{subsec:book_library}

% 400 children's books; adjective tags (atmosphere) and noun tags
% (themes) -> two LLM text embeddings -> two bundles:
% (u, v_atm, v_theme). Tag-to-description recipe
% ("This item has the following tags: ...") を setup に一文で再掲。
% Ver.2 (paper202605) の結果を移植して書き直す。
